% Options for packages loaded elsewhere
\PassOptionsToPackage{unicode}{hyperref}
\PassOptionsToPackage{hyphens}{url}
\documentclass[
  a4paper,
]{article}
\usepackage{xcolor}
\usepackage[margin=25mm]{geometry}
\usepackage{amsmath,amssymb}
\setcounter{secnumdepth}{-\maxdimen} % remove section numbering
\usepackage{iftex}
\ifPDFTeX
  \usepackage[T1]{fontenc}
  \usepackage[utf8]{inputenc}
  \usepackage{textcomp} % provide euro and other symbols
\else % if luatex or xetex
  \usepackage{unicode-math} % this also loads fontspec
  \defaultfontfeatures{Scale=MatchLowercase}
  \defaultfontfeatures[\rmfamily]{Ligatures=TeX,Scale=1}
\fi
\usepackage{lmodern}
\ifPDFTeX\else
  % xetex/luatex font selection
\fi
% Use upquote if available, for straight quotes in verbatim environments
\IfFileExists{upquote.sty}{\usepackage{upquote}}{}
\IfFileExists{microtype.sty}{% use microtype if available
  \usepackage[]{microtype}
  \UseMicrotypeSet[protrusion]{basicmath} % disable protrusion for tt fonts
}{}
\makeatletter
\@ifundefined{KOMAClassName}{% if non-KOMA class
  \IfFileExists{parskip.sty}{%
    \usepackage{parskip}
  }{% else
    \setlength{\parindent}{0pt}
    \setlength{\parskip}{6pt plus 2pt minus 1pt}}
}{% if KOMA class
  \KOMAoptions{parskip=half}}
\makeatother
\usepackage{longtable,booktabs,array}
\newcounter{none} % for unnumbered tables
\usepackage{calc} % for calculating minipage widths
% Correct order of tables after \paragraph or \subparagraph
\usepackage{etoolbox}
\makeatletter
\patchcmd\longtable{\par}{\if@noskipsec\mbox{}\fi\par}{}{}
\makeatother
% Allow footnotes in longtable head/foot
\IfFileExists{footnotehyper.sty}{\usepackage{footnotehyper}}{\usepackage{footnote}}
\makesavenoteenv{longtable}
\setlength{\emergencystretch}{3em} % prevent overfull lines
\providecommand{\tightlist}{%
  \setlength{\itemsep}{0pt}\setlength{\parskip}{0pt}}
\usepackage{bookmark}
\IfFileExists{xurl.sty}{\usepackage{xurl}}{} % add URL line breaks if available
\urlstyle{same}
\hypersetup{
  pdftitle={Emergent Structures as Factorization of an Anonymous Complex Wave Multiset: Working Hypotheses and a Numerical Verification Program for Simplices{,} Localization{,} Condensates{,} and Dynamics},
  pdfauthor={Noriaki Kihara (Independent Researcher), ORCID 0009-0004-6753-4020},
  hidelinks,
  pdfcreator={LaTeX via pandoc}}

\title{Emergent Structures as Factorization of an Anonymous Complex Wave
Multiset: Working Hypotheses and a Numerical Verification Program for
Simplices, Localization, Condensates, and Dynamics}
\author{Noriaki Kihara (Independent Researcher), ORCID
0009-0004-6753-4020}
\date{September 2, 2026 --- v1.0 --- Version DOI
10.5281/zenodo.22250744}

\begin{document}
\maketitle

\subsection{Abstract}\label{abstract}

This paper formulates a working hypothesis that asks whether structures
such as space, time, particles, vertices, edges, and complete networks
can be uniquely read out from a multiset of anonymous complex waves
alone, without being assumed as foundational.

Conventionally, when describing relations among multiple elements, one
first fixes the number of elements and counts all pairwise relations as
the number of edges of a complete graph. This description, however,
implicitly introduces identifiable and persistent vertices, endpoints
for each relation, complete connectivity, and a one-relation-one-edge
correspondence. If the underlying waves are anonymous, these cannot be
axioms; they must be reconstruction results that are admitted only after
a successful factorization of the anonymous set.

We define the base state as a complex wave multiset identified under
permutations. A submultiset is read as a harmonic series, a
localization, a condensate, a simplex, or a persistent object only when
it is uniquely factorized while simultaneously satisfying conditions
such as zero-square closure, finite recurrence, integer ratios, common
phase, low rank, and complex Gram consistency. The complete network is
not an input; it is a connectivity structure that becomes known as a
consequence, once anonymous relational quantities have been uniquely
reconstructed into a simplex.

Furthermore, we define dynamics not as trajectories of named
constituents, but as a process in which the factorization type of the
anonymous wave set is preserved while its common phase center changes.
We propose ground-truth embedding experiments in which all provenance
information of data generated from known structures is completely
erased, blind control experiments with no planted structure,
permutation-invariance tests, perturbation tests, and temporal
persistence tests. This paper does not prove the existence or uniqueness
of emergent structures. It fixes a falsifiable problem setting for
verifying them without breaking anonymity.

\textbf{Keywords}: anonymous complex waves, factorization, zero-square
closure, finite recurrence, simplex, unlabeled distance geometry,
localization, condensate, emergent dynamics, numerical experiments

\begin{center}\rule{0.5\linewidth}{0.5pt}\end{center}

\subsection{1. Introduction}\label{introduction}

When constructing a physical model, it is natural to number particles or
nodes and to describe the relations among them as arrays or matrices. If
the number of elements is N and one relation is placed between every
pair of elements, the number of relations is

\[
M=\frac{N(N-1)}{2}.
\]

This counting itself is correct. Before it, however, at least the
following assumptions have been made.

\begin{enumerate}
\def\labelenumi{\arabic{enumi}.}
\tightlist
\item
  N nodes exist in advance.
\item
  The nodes can be distinguished from one another.
\item
  The identity of each node is preserved throughout any change.
\item
  A relation exists between every pair of nodes.
\item
  Each relation has one pair of endpoints.
\item
  Relations among three or more elements are composed of pairwise
  relations.
\end{enumerate}

If, however, the only underlying entities are anonymous complex waves,
then the waves carry no vertex numbers and no edge numbers from the
start. Array indices are storage locations in a computation, not
physical names. Nevertheless, if relations are represented by index
pairs from the outset, the complete network and its connectivity
information are smuggled into the foundation.

This paper reverses that order. Instead of starting from

\[
\text{complete network}
\longrightarrow
\text{relational quantities},
\]

we ask about

\[
\text{anonymous complex wave set}
\longrightarrow
\text{unique factorization}
\longrightarrow
\text{readout of geometry, localization, condensates, and dynamics}.
\]

With this change, the complete network ceases to be an ontological
premise. Only when a part of the anonymous relational quantities is
uniquely reconstructed into a simplex with some number of vertices do we
judge that this subset could be read as the edge set of a complete
graph.

The same idea applies to harmonics. We do not first posit a base wave
and its child harmonics. When a subset of the anonymous wave set is
found that can be factorized with a common fundamental period and
integer ratios, we read it, as a result, as having been a harmonic
series. Localization, condensation, parent-child relations, and motion
are likewise defined as stable factorization results of the anonymous
set.

\begin{center}\rule{0.5\linewidth}{0.5pt}\end{center}

\subsection{2. Position and Scope of
Claims}\label{position-and-scope-of-claims}

This paper is neither a proof paper nor a report of numerical results.
It has three aims.

\begin{enumerate}
\def\labelenumi{\arabic{enumi}.}
\tightlist
\item
  Define a base state that preserves anonymity.
\item
  Define emergent structures as factorizations of the anonymous set.
\item
  Fix experimental procedures that verify the hypotheses without
  circularity and can detect failure.
\end{enumerate}

We do not claim that space, time, or particles have already been
derived. We do not claim that any particular update rule is the correct
physical law. Nor do we claim any correspondence with real particles,
the Standard Model, gravity, or cosmology.

The central proposition proposed in this paper is the following.

\begin{quote}
If a small number of strongly consistent factors exist within an
anonymous complex wave multiset, they may be uniquely read out,
independently of labels, as simplices, harmonics, localizations,
condensates, or persistent objects.
\end{quote}

\begin{center}\rule{0.5\linewidth}{0.5pt}\end{center}

\subsection{3. Anonymous Complex Wave
States}\label{anonymous-complex-wave-states}

\subsubsection{3.1 Base set}\label{base-set}

Computationally we handle M complex waves arranged in order:

\[
W=(w_1,w_2,\ldots,w_M),
\qquad
w_m\in\mathbb C.
\]

The indices, however, are not physical names. We regard any permutation
as the same state and define the base state as the following equivalence
class:

\[
\mathcal W
=
[W]_{S_M}
=
\{(w_{\pi(1)},\ldots,w_{\pi(M)})\mid \pi\in S_M\}.
\]

The physical input is therefore not an ordered vector but a multiset:

\[
\mathcal W
=
\{w_1,\ldots,w_M\}/S_M.
\]

Since several waves may share the same value, we treat this as a
multiset rather than an ordinary set.

\subsubsection{3.2 Zero-square closure}\label{zero-square-closure}

The base constraint adopted in this paper is the closure of the complex
sum of squares, not of the Hermitian norm:

\[
Q(\mathcal W)
=
\sum_{m=1}^{M}w_m^2
=0.
\]

This expression does not depend on the ordering of the waves. It is also
closed under a common rescaling by any nonzero complex number:

\[
Q(\lambda\mathcal W)
=
\lambda^2Q(\mathcal W)
=0.
\]

Hence the absolute scale is not fixed by this condition alone.

\subsubsection{3.3 Finite recurrence}\label{finite-recurrence}

We define finite recurrence not as the condition that each named wave
returns to the same place, but as the condition that the anonymous set
as a whole returns to the original state up to permutation equivalence:

\[
U^q\mathcal W
\simeq
\mathcal W.
\]

Here q is the finite recurrence order, and the equivalence includes at
least permutations of the waves. When necessary, common phase, overall
scale, or geometric congruence transformations are also included in the
equivalence relation.

Under this definition, each array element need not return to the same
index after one period. It suffices that the same state is reconstituted
as a set.

\subsubsection{3.4 Distinguishing computational order from emergent
time}\label{distinguishing-computational-order-from-emergent-time}

Numerical experiments require step numbers for performing updates:

\[
s=0,1,2,\ldots.
\]

This s, however, is the order in which the computer processed the
updates, not physical time. We call it the computational order. If, in
the future, a factor readable as time is constructed, it must be
reconstructed separately from finite cycles or reference phases within
the anonymous wave set:

\[
\tau_{\mathrm{read}}
=
F_{\mathrm{clock}}(\mathcal W_s,\mathcal W_{s+1},\ldots).
\]

\begin{center}\rule{0.5\linewidth}{0.5pt}\end{center}

\subsection{4. Anonymous Factorization}\label{anonymous-factorization}

\subsubsection{4.1 Definition of a factor}\label{definition-of-a-factor}

We express the operation of reading structure from an anonymous wave
state as a factorization map:

\[
\mathcal F:
\mathcal W
\longmapsto
\{F_1,F_2,\ldots,F_L\}.
\]

A factor F consists of a submultiset, latent parameters, a
reconstruction rule, and a residual:

\[
F
=
(C,\theta,R,\varepsilon),
\qquad
C\subseteq\mathcal W.
\]

Here C is the anonymous submultiset constituting the factor; \(\theta\)
denotes latent parameters such as a common period, a phase center,
integer orders, or a geometric dimension; R is the rule that
reconstructs C from them; and \(\varepsilon\) is the reconstruction
residual.

For an exact factor,

\[
\varepsilon(F)=0.
\]

In numerical experiments, allowing for finite precision, a factor is
called approximate when its residual is below a tolerance fixed in
advance:

\[
\varepsilon(F)\leq\varepsilon_{\mathrm{tol}}.
\]

\subsubsection{4.2 Equivalent
factorizations}\label{equivalent-factorizations}

If two factorizations differ only by permutations of the input waves,
permutations of factor names, a global phase, a common scale, a mirror
reflection, or a congruence transformation of the adopted geometry, we
regard them as the same physical interpretation.

The number of inequivalent valid factorizations is denoted by

\[
\Omega_{\mathrm{fact}}(\mathcal W)
=
\#\left(
\frac{\{\text{valid factorizations}\}}
{\text{admissible symmetries}}
\right).
\]

The judgment falls into three cases:

\[
\Omega_{\mathrm{fact}}=0
\quad\Longrightarrow\quad
\text{not reconstructible as the specified structure},
\]

\[
\Omega_{\mathrm{fact}}=1
\quad\Longrightarrow\quad
\text{uniquely reconstructed up to physical equivalence},
\]

\[
\Omega_{\mathrm{fact}}>1
\quad\Longrightarrow\quad
\text{multiple inequivalent structural interpretations exist}.
\]

Finding one planted ground truth demonstrates only existence. To claim
unique emergence, one must confirm that no inequivalent alternative
solutions exist.

\begin{center}\rule{0.5\linewidth}{0.5pt}\end{center}

\subsection{5. Simplex Reconstruction without Inputting a Complete
Network}\label{simplex-reconstruction-without-inputting-a-complete-network}

\subsubsection{5.1 Anonymous set of complex squared
lengths}\label{anonymous-set-of-complex-squared-lengths}

Take a submultiset as candidate complex squared lengths:

\[
D
=
\{d_1^2,d_2^2,\ldots,d_r^2\}.
\]

At this stage, no vertices and no edges exist, and no correspondence
between lengths and point pairs is given.

If a candidate vertex number n exists, the required number of lengths
satisfies

\[
r
=
\frac{n(n-1)}{2}.
\]

Here this equation is not a foundational axiom. It is a necessary
condition for judging whether the anonymous subset can serve as
candidates for all edges of an n-vertex simplex.

\subsubsection{5.2 Provisional edge
assignment}\label{provisional-edge-assignment}

The reconstructor examines all essentially distinct assignments from the
anonymous length set to candidate edges:

\[
f:
D
\longrightarrow
\bigl\{\{i,j\}\mid 0\leq i<j<n\bigr\}.
\]

This map is not base data; it is a candidate variable of the inverse
problem.

Taking vertex 0 as the reference, a complex Gram matrix is constructed
from the assignment f:

\[
G_{ab}
=
\frac{1}{2}
\left(
d_{0a}^2+d_{0b}^2-d_{ab}^2
\right),
\qquad
1\leq a,b<n.
\]

Only assignments satisfying the required embedding dimension, rank,
nondegeneracy, closure conditions, and any additional symmetry
conditions are accepted as valid. Because complex squared lengths are
involved, positivity and triangle inequalities valid for real Euclidean
distances are not required unconditionally. Which congruence
transformations are identified is specified in the experimental protocol
together with the adopted complex bilinear form.

\subsubsection{5.3 The complete network is a reconstruction
certificate}\label{the-complete-network-is-a-reconstruction-certificate}

Only when exactly one valid assignment remains, up to vertex
permutations and congruence transformations, do we judge that D uniquely
determines an n-vertex simplex:

\[
D
\xrightarrow{\text{unique reconstruction}}
\text{n-vertex simplex}.
\]

Only then does it become known that the elements of D corresponded to
all edges of a complete graph:

\[
K_n
\quad\text{is a reconstruction result, not an input}.
\]

The problem of reconstructing real point configurations from unlabeled
distance multisets has been studied as unlabeled distance geometry and
unlabeled global rigidity. This paper shares that problem awareness, but
adds, as working hypotheses, that the input is not restricted to
positive real distances and that complex waves, zero-square closure,
finite recurrence, harmonic factors, and factor persistence are treated
within a single readout framework.

\begin{center}\rule{0.5\linewidth}{0.5pt}\end{center}

\subsection{6. Factorization of Harmonics, Localization, and
Condensation}\label{factorization-of-harmonics-localization-and-condensation}

\subsubsection{6.1 Harmonics are not an a priori parent-child
relation}\label{harmonics-are-not-an-a-priori-parent-child-relation}

We do not posit a base wave first and generate integer multiples as
child waves. From the phase histories or cyclic properties of the
anonymous wave set, we search for a subset C for which a common
fundamental quantity and an integer sequence exist:

\[
\omega_m
\simeq
k_m\omega_0,
\qquad
k_m\in\mathbb Z,
\qquad
w_m\in C.
\]

A wave corresponding to the fundamental quantity need not be explicitly
present in the set. The fundamental can be inferred inversely as a
latent factor shared by several waves.

For phases too, we search for the following low-dimensional
representation:

\[
\phi_m(s)
\simeq
k_m\phi_0(s)+\delta_{k_m}
\pmod{2\pi}.
\]

If this representation holds with a small residual, C is read as a
harmonic factor.

\subsubsection{6.2 Localization is a readout of common
phase}\label{localization-is-a-readout-of-common-phase}

The readout kernel obtained by synthesizing a harmonic factor is written
as

\[
L_C(\phi)
=
\left|
\sum_{m\in C}
A_m
\exp\left[
i k_m(\phi-\phi_0)
+i\delta_{k_m}
\right]
\right|^2.
\]

When L has a sharp peak near the phase \(\phi_0\), the factor is read as
a localized state. This does not mean that the underlying waves exist at
a point in a background space. It means that the common phase factor of
the anonymous wave set has a sharp coincidence position with respect to
a particular readout.

If the highest harmonic order is K, then for an ideal series with equal
amplitudes and regular phases, the phase resolution improves roughly at
the order

\[
\Delta\phi
\sim
\frac{1}{K}.
\]

Even when many harmonics are present, if they are all generated from one
common fundamental phase and a fixed integer sequence, the number of
independent degrees of freedom need not have grown with the number of
harmonics. The very fact that many components can be described by a few
latent factors serves as an indicator of condensation or low-rankness.

\subsubsection{6.3 Condensates}\label{condensates}

In this paper a condensate is defined not as a fixed set of named
elements but as a persistent low-dimensional factor satisfying the
following.

\begin{enumerate}
\def\labelenumi{\arabic{enumi}.}
\tightlist
\item
  Many anonymous waves can be reconstructed from a few latent
  parameters.
\item
  The partial zero-square closure, or a specified conservation residual,
  is small.
\item
  It has a common phase center and a finite phase width.
\item
  Its factor signature, modulo admissible symmetries, is preserved
  across computational order.
\end{enumerate}

The constituent waves of a condensate may be exchanged during updates.
What is conserved is not the names of elements but the factorization
type.

\begin{center}\rule{0.5\linewidth}{0.5pt}\end{center}

\subsection{7. Symmetry and Uniqueness}\label{symmetry-and-uniqueness}

\subsubsection{7.1 The need for strong
symmetry}\label{the-need-for-strong-symmetry}

Sets that simultaneously satisfy zero-square closure, finite recurrence,
integer ratios, low rank, and geometric consistency are expected to be
special among generic random complex wave sets. Strong symmetry is a
candidate condition making it possible to compress many anonymous
components into a few factors.

On the other hand, perfect symmetry produces a degeneracy in which no
vertex, direction, or phase center can be distinguished. What is needed
is therefore not simply maximal symmetry. We take as a working
hypothesis the two-layer structure

\[
\text{strong symmetry enabling factorization}
+
\text{a small symmetry breaking selecting the readout}.
\]

\subsubsection{7.2 Identification of symmetric
solutions}\label{identification-of-symmetric-solutions}

The many solutions obtained by renaming the vertices of a regular
simplex are not different physical structures. Uniqueness is therefore
judged not by the number of labeled solutions but by the number of
equivalence classes after quotienting by admissible symmetries.

If, however, distinct simplices, distinct harmonic fundamentals, or
distinct condensate partitions remain that cannot be connected by
symmetry transformations, that is genuine non-uniqueness.

\subsubsection{7.3 Perturbation regimes}\label{perturbation-regimes}

Starting from a perfectly symmetric state, we add perturbations that
preserve anonymity:

\[
w_m(\epsilon)
=
w_m^{(0)}+\epsilon\eta_m.
\]

We test for the possible existence of three regimes.

\begin{enumerate}
\def\labelenumi{\arabic{enumi}.}
\tightlist
\item
  Perfectly symmetric point: the shape is clear but directions or
  centers are degenerate.
\item
  Small-perturbation regime: the shape is preserved while a position or
  direction is selected.
\item
  Large-perturbation regime: valid factors are lost and reconstruction
  becomes non-unique or impossible.
\end{enumerate}

The existence of a critical perturbation magnitude is itself a working
hypothesis and may be refuted by experiment.

\begin{center}\rule{0.5\linewidth}{0.5pt}\end{center}

\subsection{8. Dynamics as Preservation of Factorization
Type}\label{dynamics-as-preservation-of-factorization-type}

\subsubsection{8.1 Microscopic updates and observed
dynamics}\label{microscopic-updates-and-observed-dynamics}

A microscopic update is expressed as a map producing a sequence of
anonymous states:

\[
\mathcal W_{s+1}
=
\Phi(\mathcal W_s).
\]

The update map must not depend on input ordering:

\[
\Phi(\pi\mathcal W)
\simeq
\Phi(\mathcal W).
\]

At each computational order we perform a factorization:

\[
\mathcal F(\mathcal W_s)
=
\{F_1(s),F_2(s),\ldots\}.
\]

Observed dynamics is defined not as trajectories of individual waves but
as changes of factor equivalence classes:

\[
[F_\alpha(s)]
\longrightarrow
[F_\beta(s+1)].
\]

\subsubsection{8.2 Factor signature}\label{factor-signature}

To judge the identity of a factor, we define a signature that does not
depend on element names:

\[
\Sigma(F)
=
\left(
q,
\{k_m\},
\text{amplitude ratios},
\text{closure residual},
\text{Gram spectrum},
\text{rank},
\text{localization width},
\ldots
\right).
\]

The concrete signature components are preregistered for each target
factor. When the signature agrees before and after an update and only
the common phase center changes, the factor is read as a conserved
condensate.

\subsubsection{8.3 Motion}\label{motion}

With factor shape \(C_0\) and common phase translation T, a moving
factor is expressed as

\[
F(s)
\simeq
T_{\phi(s)}F^{(0)}.
\]

If after the update

\[
\Sigma(F(s+1))
\simeq
\Sigma(F(s))
\]

still holds, the lump is conserved. Meanwhile the center phase may
change:

\[
\phi(s+1)
\neq
\phi(s).
\]

Motion is therefore not the same named material changing its position;
it is the persistence of the same factorization type while its common
phase center changes.

\subsubsection{8.4 Fusion, splitting, and
interactions}\label{fusion-splitting-and-interactions}

Recompositions of the factorization are read as follows:

\[
\{F_A,F_B\}
\longrightarrow
\{F_C\}
\qquad
\text{fusion candidate},
\]

\[
\{F_A\}
\longrightarrow
\{F_B,F_C\}
\qquad
\text{splitting candidate},
\]

\[
\{F_A,F_B\}
\longrightarrow
\{F'_A,F'_B\}
\qquad
\text{scattering or exchange candidate}.
\]

To call these physical interactions, one must identify global invariants
that are conserved independently of changes in the number of factors.
This paper does not assume their concrete form.

\begin{center}\rule{0.5\linewidth}{0.5pt}\end{center}

\subsection{9. Working Hypotheses}\label{working-hypotheses}

We divide the claims of this paper into forms that future experiments
can refute individually.

\subsubsection{Hypothesis H1: Anonymous simplex
reconstruction}\label{hypothesis-h1-anonymous-simplex-reconstruction}

Some anonymous sets of complex squared quantities with strong
consistency are uniquely reconstructed into a single complex simplex, up
to vertex permutations, congruence transformations, and admissible
gauges.

\subsubsection{Hypothesis H2: A posteriori appearance of the complete
network}\label{hypothesis-h2-a-posteriori-appearance-of-the-complete-network}

For a subset admitting a unique simplex reconstruction, the complete
graph structure becomes known as a result, from the number of elements
and the connection assignment. There is no need to input a complete
network.

\subsubsection{Hypothesis H3: Localization readout by harmonic
factors}\label{hypothesis-h3-localization-readout-by-harmonic-factors}

An anonymous wave subset sharing integer ratios, a common fundamental
phase, and regular amplitudes produces a phase peak of finite width in
the synthesized readout. The center and width of this peak are read as
the localization position and localization width.

\subsubsection{Hypothesis H4: Condensates as low-rank
factors}\label{hypothesis-h4-condensates-as-low-rank-factors}

When many anonymous waves are reconstructible from a few latent factors
and the closure and phase width remain stable, the factor persists as a
condensate.

\subsubsection{Hypothesis H5: Motion as preservation of factor
type}\label{hypothesis-h5-motion-as-preservation-of-factor-type}

When the signature of a condensed factor is preserved and only its
common phase center changes continuously, the change is read as motion.

\subsubsection{Hypothesis H6: Emergent space, time, and
objects}\label{hypothesis-h6-emergent-space-time-and-objects}

When a simplex factor, an internal cyclic factor, and a persistent
condensed factor are independently reconstructed from the same anonymous
wave sequence and are mutually consistent, they can be read respectively
as space, time, and objects.

\subsubsection{Hypothesis H7: Symmetry with small
breaking}\label{hypothesis-h7-symmetry-with-small-breaking}

There exists a finite regime in which strong symmetry enabling unique
factorization coexists with a small symmetry breaking that selects
positions and directions.

\begin{center}\rule{0.5\linewidth}{0.5pt}\end{center}

\subsection{10. Numerical Verification
Program}\label{numerical-verification-program}

\subsubsection{10.1 Complete separation of generator and
reconstructor}\label{complete-separation-of-generator-and-reconstructor}

The experiments consist of two independent programs.

The generator produces complex wave sets from known structures:

\[
S_{\mathrm{true}}
\xrightarrow{\mathrm{generator}}
\mathcal W.
\]

After generation, all of the following are erased:

\begin{itemize}
\tightlist
\item
  vertex numbers
\item
  edge endpoints
\item
  generation order of waves
\item
  harmonic orders
\item
  fundamental-wave labels
\item
  condensate labels
\item
  parent-child labels
\item
  individual IDs across computational order
\item
  coordinates used in generation
\end{itemize}

The input ordering is further randomized.

The reconstructor receives only the permitted complex values, phase
histories, and global constraints. The reconstructor must not access the
generator's internal data, ground-truth labels, or connection tables.

\subsubsection{10.2 Experiment E0: Permutation-invariance
test}\label{experiment-e0-permutation-invariance-test}

The same anonymous set is input under multiple random permutations:

\[
\mathcal W^{(r)}
=
\{w_{\pi_r(1)},\ldots,w_{\pi_r(M)}\}.
\]

We require that all inputs yield the same factorization result up to
admissible symmetries. If inequivalent results depend on the input
ordering, the implementation has a hidden label dependence.

\subsubsection{10.3 Experiment E1: Ground-truth embedded
simplex}\label{experiment-e1-ground-truth-embedded-simplex}

All pairwise relational quantities are generated from a known real or
complex simplex, and the correspondence table is erased. We test whether
the reconstructor satisfies the following.

\begin{enumerate}
\def\labelenumi{\arabic{enumi}.}
\tightlist
\item
  It recovers a solution congruent to the original simplex.
\item
  No inequivalent alternative solution exists.
\item
  Even without inputting the vertex number, it recovers the candidate
  vertex number.
\item
  The complete network is obtained as an output result.
\end{enumerate}

\subsubsection{10.4 Experiment E2: Ground-truth embedded harmonics and
localization}\label{experiment-e2-ground-truth-embedded-harmonics-and-localization}

Wave sets are generated from a known fundamental phase and integer
orders and then anonymized. We test whether the reconstructor can
recover the fundamental quantity, the set of integer orders, the phase
center, and the localization width.

We also test the case where the wave corresponding to the fundamental is
removed from the input set, checking whether the latent fundamental
factor can be inferred inversely.

\subsubsection{10.5 Experiment E3: Composite
factors}\label{experiment-e3-composite-factors}

A simplex factor, a harmonic factor, and a closure factor are
simultaneously embedded in the same anonymous wave set. We verify
whether the individually recovered factors occupy the same subsets or
consistent overlaps.

This experiment examines whether different readouts point to the same
emergent structure.

\subsubsection{10.6 Experiment E4: Persistence and motion
lock}\label{experiment-e4-persistence-and-motion-lock}

We generate sequences of anonymous sets in which the factor shape is
preserved and only the common phase center is varied. At each
computational order, the wave ordering and the constituents are
shuffled, making named tracking impossible.

We verify whether the reconstructor can recover the same factor
signature and the change of phase center without using the identity of
the constituent waves.

\subsubsection{10.7 Experiment E5: Blind
controls}\label{experiment-e5-blind-controls}

Without planting any structure, we generate anonymous sets satisfying
only zero-square closure, finite recurrence, or the same marginal
distributions. Compared with the ground-truth embedded group, we measure
the frequencies of valid factors, unique factors, and false positives:

\[
P_{\mathrm{unique}}
=
P\left(
\Omega_{\mathrm{fact}}=1
\mid
\text{generating condition}
\right).
\]

If structures are detected at the same frequency in the control group as
in the ground-truth group, the reconstruction conditions are too weak.

\subsubsection{10.8 Experiment E6: Perturbation and
criticality}\label{experiment-e6-perturbation-and-criticality}

Controlled perturbations are added to ground-truth embedded sets, and
the following are measured:

\begin{itemize}
\tightlist
\item
  number of factorizations
\item
  ground-truth recovery rate
\item
  reconstruction residual
\item
  Gram rank and spectrum
\item
  localization peak width
\item
  lifetime of factor signatures
\item
  false-positive rate
\end{itemize}

The regime from the perfectly symmetric point to factor collapse is
scanned continuously.

\begin{center}\rule{0.5\linewidth}{0.5pt}\end{center}

\subsection{11. Success Conditions}\label{success-conditions}

The minimal success conditions of this research program are set as
follows.

\begin{enumerate}
\def\labelenumi{\arabic{enumi}.}
\tightlist
\item
  Changing the input ordering yields the same physical equivalence
  class.
\item
  The ground-truth structure can be recovered without using names,
  coordinates, or connectivity information from generation.
\item
  No alternative solution exists that is not symmetry-equivalent to the
  ground truth.
\item
  Multiple independent readouts point to the same factor.
\item
  Factors have a finite stability regime against small perturbations.
\item
  Persistent objects can be identified solely by conservation of factor
  signatures, without tracking constituent waves.
\item
  The computational step itself is not taken as physical time; a time
  candidate can be reconstructed from internal cycles.
\item
  Verifiable predictions about the next factor state can be obtained
  from the current factor state.
\end{enumerate}

Constructing separate readouts for space, time, and objects is not
sufficient. We require that they be mutually consistent on the same
anonymous wave set and constrain each other's predictions.

\begin{center}\rule{0.5\linewidth}{0.5pt}\end{center}

\subsection{12. Refutation Conditions}\label{refutation-conditions}

If any of the following holds, the corresponding hypothesis is not
supported.

\begin{enumerate}
\def\labelenumi{\arabic{enumi}.}
\tightlist
\item
  Inequivalent structures emerge depending on the ordering of the input
  array.
\item
  Even when the ground truth is recovered, many equally valid
  inequivalent solutions remain.
\item
  Structures can be recovered only if ground-truth labels or connection
  tables are passed to the reconstructor.
\item
  Structures are detected in random control groups at rates comparable
  to ground-truth embedded groups.
\item
  Factors are completely destroyed by arbitrarily small perturbations.
\item
  The localization readout depends on external coordinates used at
  generation time.
\item
  The judgment of motion depends on tracking named constituent waves.
\item
  Emergent time is nothing but an alias of the computational step.
\item
  The simplex readout, the harmonic localization readout, and the
  persistent factor readout contradict one another.
\item
  Expanding the search range destroys uniqueness, revealing that the
  earlier unique solution was an artifact of insufficient search.
\end{enumerate}

Failure results are also important. The conditions under which
factorization becomes impossible, non-unique, or unstable delimit the
regime in which emergent structures hold.

\begin{center}\rule{0.5\linewidth}{0.5pt}\end{center}

\subsection{13. Computational Cost and Search
Strategy}\label{computational-cost-and-search-strategy}

A naive search assigning r anonymous candidate lengths to the edges of
an n-vertex simplex grows rapidly. In the generic case where all values
are distinct, an estimate of the number of candidates up to vertex
permutations is

\[
\frac{r!}{n!},
\qquad
r=\frac{n(n-1)}{2}.
\]

For large systems, therefore, instead of exhaustive permutation search,
the following pruning becomes necessary:

\begin{itemize}
\tightlist
\item
  partial Gram consistency
\item
  rank conditions
\item
  subsets with zero-square closure
\item
  finite cycle orders
\item
  integer-ratio candidates
\item
  repeated values and automorphisms
\item
  spectral invariants
\item
  low-rank decompositions
\item
  gluing consistency of partial simplices
\end{itemize}

In ground-truth embedding experiments the generator knows the answer,
but the reconstructor must not use that answer for pruning. The ground
truth is used only for evaluation after the reconstruction has finished.

When the search becomes large-scale, constructing uniqueness
certificates instead of enumerating all solutions will also be
considered. However, uniqueness must not be claimed merely because one
locally optimal solution has been found.

\begin{center}\rule{0.5\linewidth}{0.5pt}\end{center}

\subsection{14. Minimal Relation to Existing
Work}\label{minimal-relation-to-existing-work}

The question of whether real point configurations can be reconstructed
from unlabeled distance multisets has been studied. Boutin and Kemper
treated the problem of reconstructing point configurations from the
distribution of distances or areas, showing that many generic
configurations are reconstructible while counterexamples also exist.
Gortler, Theran, and Thurston studied the relation between unlabeled
distances and global rigidity in generic position.

This paper does not claim that these results apply as they stand to
complex wave systems. Handling complex squared quantities lacking the
positivity of real positive distances, zero-square closure, finite
recurrence, harmonic factors, and factor persistence across time
simultaneously requires a different formulation and numerical
verification.

\begin{center}\rule{0.5\linewidth}{0.5pt}\end{center}

\subsection{15. Limitations and Future
Work}\label{limitations-and-future-work}

At the stage of this paper, the following remain unresolved.

\begin{enumerate}
\def\labelenumi{\arabic{enumi}.}
\tightlist
\item
  Whether anonymity should quotient only by permutations, or whether
  wider basis transformations should also be identified.
\item
  The congruence group of complex squared lengths and the exact
  equivalence relation for unique reconstruction.
\item
  Whether the global structure in the case of overlapping factors
  becomes a partition, a covering, or a hypergraph.
\item
  Natural measures and tolerances for valid factor residuals.
\item
  Conditions under which strong symmetry increases unique
  factorizations, versus conditions under which it increases degeneracy.
\item
  Concrete procedures for constructing emergent time from internal
  cycles.
\item
  Microscopic update rules that preserve anonymity while producing
  interactions between factors.
\item
  Global invariants conserved even when the number of factors changes.
\item
  Uniqueness-certifying algorithms enabling large-scale search.
\end{enumerate}

These are not concealments of gaps in this paper; they are explicit
tasks connecting the working hypotheses to the next experiments.

\begin{center}\rule{0.5\linewidth}{0.5pt}\end{center}

\subsection{16. Conclusion}\label{conclusion}

This paper proposed a research framework that does not place complete
networks, harmonic parent-child relations, localized bodies, or
particles at the foundation, but reads them out solely from an anonymous
complex wave multiset.

Taking

\[
\mathcal W
=
\{w_1,\ldots,w_M\}/S_M
\]

as the base state, valid partial factors are searched for:

\[
\mathcal W
\xrightarrow{\text{anonymous factorization}}
\left\{
\begin{array}{l}
\text{simplex factors}\\
\text{harmonic factors}\\
\text{localization factors}\\
\text{condensate factors}\\
\text{persistent factors}
\end{array}
\right\}.
\]

When a simplex factor is uniquely reconstructed, the complete network
turns out to be its reconstruction result, not an input. Harmonic
relations likewise appear, without positing a parent-child structure
first, as factorization results of subsets sharing a common fundamental
quantity and integer ratios. Localization is a sharp readout of a common
phase factor, and a condensate is a low-rank factor describing many
anonymous waves by a few latent quantities.

Dynamics is defined not as trajectories of named waves but as changes of
factor equivalence classes that preserve the factor signature while
changing the common phase center:

\[
[F(s)]
\longrightarrow
[F(s+1)].
\]

The success or failure of this framework can be judged by blind
reconstruction of anonymized ground-truth data, control experiments with
no planted structure, permutation-invariance tests, perturbation tests,
and persistence tests.

The central proposal of this paper is not to assume structures and
assign them to waves. It is to ask

\[
\boxed{
\begin{array}{l}
\text{under which conditions an anonymous wave set}\\
\text{is uniquely factorized as a single structure.}
\end{array}
}
\]

If factors readable as space, time, and particles are reconstructed
consistently with one another from the same anonymous wave set, without
using names or coordinates from generation time, the minimal goal of
this research program is achieved.

\begin{center}\rule{0.5\linewidth}{0.5pt}\end{center}

\subsection{References}\label{references}

\begin{enumerate}
\def\labelenumi{\arabic{enumi}.}
\item
  Mireille Boutin and Gregor Kemper, ``On Reconstructing n-Point
  Configurations from the Distribution of Distances or Areas,''
  \emph{Advances in Applied Mathematics}, Vol. 32, No.~4, pp.~709--735,
  2004. \url{https://doi.org/10.1016/S0196-8858(03)00101-5}
\item
  Steven J. Gortler, Louis Theran, and Dylan P. Thurston, ``Generic
  Unlabeled Global Rigidity,'' \emph{Forum of Mathematics, Sigma}, Vol.
  7, e21, 2019. \url{https://doi.org/10.1017/fms.2019.16}
\end{enumerate}

\begin{center}\rule{0.5\linewidth}{0.5pt}\end{center}

\subsection{Appendix A: List of
Symbols}\label{appendix-a-list-of-symbols}

{\def\LTcaptype{none} % do not increment counter
\begin{longtable}[]{@{}
  >{\raggedright\arraybackslash}p{(\linewidth - 2\tabcolsep) * \real{0.5000}}
  >{\raggedright\arraybackslash}p{(\linewidth - 2\tabcolsep) * \real{0.5000}}@{}}
\toprule\noalign{}
\begin{minipage}[b]{\linewidth}\raggedright
Symbol
\end{minipage} & \begin{minipage}[b]{\linewidth}\raggedright
Meaning
\end{minipage} \\
\midrule\noalign{}
\endhead
\bottomrule\noalign{}
\endlastfoot
M & total number of anonymous complex waves in the base input \\
W & ordered complex wave array in computation \\
\(\mathcal W\) & anonymous complex wave multiset identified under
permutations \\
Q & closure quantity given by the complex sum of squares \\
q & finite recurrence order \\
s & update order in numerical computation; not physical time \\
F & factor extracted from the anonymous set \\
C & submultiset constituting a factor \\
\(\theta\) & latent parameters of a factor \\
\(\varepsilon\) & factor reconstruction residual \\
\(\Omega_{\mathrm{fact}}\) & number of inequivalent factorizations
modulo admissible symmetries \\
D & anonymous set of complex squared quantities used for simplex
reconstruction \\
n & candidate number of simplex vertices for reconstruction \\
G & complex Gram matrix built from a candidate assignment \\
k & integer order of a harmonic factor \\
\(\phi_0\) & recovered common phase center \\
\(\Sigma\) & invariant signature of a factor \\
\end{longtable}
}

\begin{center}\rule{0.5\linewidth}{0.5pt}\end{center}

\subsection{Appendix B: Items to Preregister for
Experiments}\label{appendix-b-items-to-preregister-for-experiments}

Before running each numerical experiment, at least the following are
fixed.

\begin{enumerate}
\def\labelenumi{\arabic{enumi}.}
\tightlist
\item
  Data items permitted as input
\item
  Generation information hidden from the reconstructor
\item
  Symmetry transformations to be identified
\item
  Definition of factor residuals
\item
  Numerical tolerances
\item
  Search ranges and cutoff conditions
\item
  Method for judging uniqueness
\item
  Method for evaluating ground-truth recovery
\item
  Method for generating random control groups
\item
  Perturbation distributions
\item
  Success conditions
\item
  Refutation conditions
\end{enumerate}

This preregistration prevents the circularity of conveniently changing
factor definitions, tolerances, or search ranges after seeing the
results.

\end{document}
